\emph{证明系统}提供了一种纯句法方法，来刻画语句之间的后承与相容关系。\emph{相继式演算}便是此类证明系统之一。其中的一个\emph{推导}是一棵由相继式组成的树（一个相继式 $\Gamma \Sequent \Delta$ 由两个公式序列经~$\Sequent$ 隔开构成）。推导树最顶端的相继式是形如 $!A \Sequent !A$ 的\emph{初始相继式}。在一个正确的推导中，其余所有相继式都必须由某条\emph{推理规则}正确地提供依据。这些规则成对出现：每个联结词和量词都有一条处理相继式左侧的规则和一条处理其右侧的规则。例如，如果一个相继式 $\Gamma \Sequent \Delta, !A \lif !B$ 是由 $\RightR{\lif}$ 规则得出的，则其上方的相继式（即\emph{前提}）必须是 $!A, \Gamma \Sequent \Delta, !B$。有些规则还允许调整相继式中语句的顺序或数量，例如，$\RightR{\Exchange}$ 规则就允许交换相继式右侧的两个公式。

如果存在相继式 $\quad \Sequent !A$ 的推导，我们就称 $!A$ 是一个\emph{定理}，并记作~$\Proves !A$。如果存在 $\Gamma_0 \Sequent !A$ 的推导，且 $\Gamma_0$ 中的每个 $!B$ 都属于~$\Gamma$，我们就称 $!A$ \emph{可从}~$\Gamma$ \emph{推导}，并记作 $\Gamma \Proves !A$。如果存在 $\Gamma_0 \Sequent \quad$ 的推导，且 $\Gamma_0$ 中的每个 $!B$ 都属于~$\Gamma$，我们就称 $\Gamma$ 是\emph{不一致的}，否则称其为\emph{一致的}。这些概念是相互联系的，例如，$\Gamma \Proves !A$ 当且仅当 $\Gamma \cup \{\lnot !A\}$ 是不一致的。它们也与对应的语义概念相关，例如，如果 $\Gamma \Proves !A$，那么 $\Gamma \Entails !A$。证明系统的这一性质——从 $\Gamma$ 可推导出的必然为 $\Gamma$ 所语义后承——被称为\emph{可靠性}。\emph{可靠性定理}是通过对推导的长度作归纳来证明的，即证明：若前提相继式有效，则每一个单独的推理都能保持结论相继式的有效性。